
%%% THIS DOCUMENT NEEDS TO BE COMPILED WITH LUATEX OR LUALATEX ENGINE

\documentclass[10pt]{amsart}
% Math Packages
\usepackage{tcolorbox}
\usepackage{amsmath, mathtools}
\usepackage{amssymb}
\usepackage{amsthm}
\usepackage{amsfonts}
\usepackage{bbm}
\usepackage{breqn}
\usepackage[margin=1in]{geometry}
\usepackage{graphicx}
\usepackage{tikz}
\usetikzlibrary{arrows.meta}
\usetikzlibrary{calc}
\usepackage{forest}
\usepackage{tikz-qtree}
\graphicspath{ {./images/} }
\usepackage{hyperref}
\usepackage[capitalize]{cleveref}
\usepackage[shortlabels]{enumitem}
\usetikzlibrary{arrows,matrix,positioning}
\usepackage{multicol}

% for the pipe symbol
\usepackage[T1]{fontenc}

% Colorful Notes
\usepackage{color} \definecolor{Red}{rgb}{1,0,0} \definecolor{Blue}{rgb}{0,0,1}
\definecolor{Purple}{rgb}{.5,0,.5} \def\red{\color{Red}} \def\blue{\color{Blue}}
\def\gray{\color{gray}} \def\purple{\color{Purple}}
\newcommand{\rnote}[1]{{\red [#1]}} % \rnote{foo} gives '[foo]' in red
\newcommand{\pnote}[1]{{\purple [#1]}} % \pnote{foo} gives '[foo]' in purple
\newcommand{\bnote}[1]{{\blue #1}} % \bnote{foo} gives 'foo' in blue
\newcommand{\gnote}[1]{{\gray #1}} % \gnote{foo} gives 'foo' in gray
\newcommand{\Max}[1]{{\purple [#1]}} % \bnote{foo} then 'foo' is blue


% Claim numbering (the counter restarts after each proof environment)
\newcounter{claimcount}
\setcounter{claimcount}{0}
\newenvironment{claim}{\refstepcounter{claimcount}\par\addvspace{\medskipamount}\noindent\textbf{Claim \arabic{claimcount}:}}{}
\usepackage{etoolbox}
\AtBeginEnvironment{proof}{\setcounter{claimcount}{0}}
\newenvironment{claimproof}{\par\addvspace{\medskipamount}\noindent\textit{Proof of Claim  \arabic{claimcount}.}}{\hfill\ensuremath{\qedsymbol} \tiny{Claim}

  \medskip}
% Add claim support to cleverref
\crefname{claimcount}{Claim}{Claims}


% Math Environments
\newtheorem{theorem}{Theorem}
\newtheorem{assumption}[theorem]{Assumption}
\newtheorem{lemma}[theorem]{Lemma}
\newtheorem{proposition}[theorem]{Proposition}
\newtheorem{corollary}[theorem]{Corollary}
\newtheorem{question}[theorem]{Question}
\theoremstyle{definition}
\newtheorem{definition}[theorem]{Definition}
\newtheorem{remark}[theorem]{Remark}
\newtheorem{example}[theorem]{Example}
\newtheorem{notation}[theorem]{Notation}
\newtheorem{problem}[theorem]{Problem}

% Matrices and Column Vectors. 
\usepackage{stackengine}
\setstackgap{L}{1.0\normalbaselineskip}
\usepackage{tabstackengine}
\setstackEOL{;}% row separator
\setstackTAB{,}% column separator
\setstacktabbedgap{1ex}% inter-column gap
\setstackgap{L}{1.5\normalbaselineskip}% inter-row baselineskip
\let\nmatrix\bracketMatrixstack  %Usage: \nmatrix{1,2,3\4,5,6}
\newcommand\cv[1]{\setstackEOL{,}\bracketMatrixstack{#1}} %usage: \cv{1,2,3}

% table alignment
\usepackage{array}
\usepackage{booktabs}

% Custom Math Coqmmands
\newcommand{\vt}{\vskip 5mm} % vertical space
\newcommand{\fl}{\noindent\textbf} % first line
\newcommand{\Fl}{\vt\noindent\textbf} % first line with space above
\newcommand{\norm}[1]{\left\lVert#1\right\rVert} % norm
\newcommand{\pnorm}[1]{\left\lVert#1\right\rVert_p} % p-norm
\newcommand{\qnorm}[1]{\left\lVert#1\right\rVert_q} % q-norm
\newcommand{\1}[1]{\textbf{1}_{\left[#1\right]}} % indicator function 
\def\limn{\lim_{n\to\infty}} % shortcut for lim as n-> infinity
\def\sumn{\sum_{n=1}^{\infty}} % shortcut for sum from n=1 to infinity
\def\sumkn{\sum_{k=1}^{n}} % shortcut for sum from k=1 to n
\def\sumin{\sum_{i=1}^{n}} % shortcut for sum from i=1 to n
\def\SAs{\sigma\text{-algebras}} % shortcut for $\sigma$-algebras
\def\SA{\sigma\text{-algebra}} % shortcut for $\sigma$-algebra
\def\Ft{\mathcal{F}_t} % time-indexed sigma-algebra (t)
\def\Fs{\mathcal{F}_s} % time-indexed sigma-algebra (s)
\def\F{\mathcal{F}} % sigma-algebra
\def\G{\mathcal{G}} % sigma-algebra
\def\R{\mathbb{R}} % Real numbers
\def\N{\mathbb{N}} % Natural numbers
\def\Z{\mathbb{Z}} % Integers
\def\E{\mathbb{E}} % Expectation
\def\P{\mathbb{P}} % Probability
\def\Q{\mathbb{Q}} % Q probability
\def\dist{\text{dist}} %Text 'dist' for things like 'dist(x,y)'
\newcommand{\indep}{\perp \!\!\! \perp}  %independence symbol
\def\Var{\mathrm{Var}} % Variance
\def\tr{\mathrm{tr}} % trace

% Brackets and Parentheses
\def\[{\left [}
    \def\]{\right ]}
% \def\({\left (}
%   \def\){\right )}



\usepackage{color}
\definecolor{Red}{rgb}{1,0,0}
\definecolor{Blue}{rgb}{0,0,1}
\definecolor{Purple}{rgb}{.5,0,.5}
\def\red{\color{Red}}
\def\blue{\color{Blue}}
\def\gray{\color{gray}}
\def\purple{\color{Purple}}
\definecolor{RoyalBlue}{cmyk}{1, 0.50, 0, 0}
\newcommand{\dempfcolor}[1]{{\color{RoyalBlue}#1}} 
\newcommand{\demph}[1]{\dempfcolor{\textbf{\textit{#1}}}}



\begin{document}
\begin{center}
  \section*{Worksheet 2: Hypothesis tests}
\end{center}

\textit{\textbf{Instructions:} This worksheet presents three examples of
  hypothesis tests. Please work in groups. You will be asked to turn in your
  group worksheet on April 10.}



\phantom{don't write in starlight, cause the words may come out real}
\begin{center}
  \begin{tcolorbox}[ width=1\textwidth, colback=white, colframe=black,
    title=Background on chi-squared statistic, colbacktitle=blue!15,
    coltitle=black ]
    Suppose we have a multinomial random variable with $k$ outcomes, labelled
    $1,\ldots,k$, constructed from $n$ trials such that in each trial outcome
    $i$ occurs with probability $p_{i}$. Recall, the \demph{chi-squared
      statistic} is defined as
    \begin{equation*}
      \chi^{2}:= \sum_{i=1}^{k} \frac{\left(N_{i}-E_{i}\right)^{2}}{E_{i}},
    \end{equation*}
    where $E_{i} = (\text{expected \# of times outcome $i$ occurs})$, and
    $N_{i}= (\text{\# of times outcome $i$ is observed})$.

    We'll illustrate the use of this statistic in two different types of
    hypothesis tests:
    \begin{itemize}
      \item For a ``goodness of fit test'', $\chi^{2}$ has a chi-squared
      distribution with $k-1$ degrees of freedom.
      \item For a ``test of independence'' using an $r\times c$ contingency
      table, $\chi^{2}$ has a chi-squared distribution with $(r-1)\times(c-1)$
      degrees of freedom.
    \end{itemize}

    Recall, the pdf of a chi-squared random variable with $d$ degrees of
    freedom is
    \begin{equation*}
      f(x) = \left\{ \begin{array}{l@{\quad:\quad}l} \frac{1}{2^{d/2}\Gamma(d/2)}x^{\frac{d}{2}-1}e^{-x/2} & x \geq 0\\ 0& x<0\end{array}\right..
    \end{equation*}
  \end{tcolorbox}
\end{center}

\subsection{Hypothesis test type \#1: A ``goodness of fit'' test}


\vspace{1em}
\begin{problem}[Goodness-of-fit test]
  According to the \textit{Mars Wrigley Confectionery Company}, the five
  colors of skittles are produced in equal proporitions. Is that really true?
  

  Let's take their claim to be the null hypothesis.
  \begin{equation*}
    H_{0}: p_{\rm green} = p_{\rm yellow}=p_{\rm red} = p_{\rm orange} = p_{\rm purple} = 0.2.
  \end{equation*}
  The alternative hypothesis is
  \begin{equation*}
    H_{a}: \text{at least one of the probabilities is different from $0.2$}
  \end{equation*}

  \begin{enumerate}[(a)]
    \setlength{\itemsep}{1em}
    \item Go grab a random sample of at least 30 skittles. What is your sample
    size?

    \item For each
    $i\in \left\{\text{green}, \text{yellow}, \text{red}, \text{orange},
      \text{purple}\right\}$, let $E_{i}$ denote the expected number of
    skittles of color $i$ in your sample, assuming that $H_{0}$ is true. Find
    $E_{i}$ for all $i$.
    
    \item Fill in the following table with expected and observed counts:
    \begin{center}
      \begin{tabular}{l|l|l|l|l|l|}
        & \textbf{green} & \textbf{yellow} & \textbf{  red  } & \textbf{orange} & \textbf{purple} \\ \hline
        \textbf{expected} &  $\phantom{\displaystyle\int}$ &   &       &   &   \\ \hline
        \textbf{observed} &  $\phantom{\displaystyle\int}$ &   &       &   &   \\ \hline
      \end{tabular}
    \end{center}
    \item Compute the chi-squared statistic $\chi^{2}$ for your sample of
    skittles. \label{item:4}
    \item Circle the correct phrase to complete the sentence: \textit{A large
      chi-square value means the observed data is (very different from / very
      similar to) the expected data.}
    \item Under the null hypothesis of a goodness-of-fit test, $\chi^{2}$ has
    a chi-squared distribution with $k-1$ degrees of freedom. Set up and
    evaluate an integral to determine the p-value
    $\P\left[\chi^{2}\geq u\right]$, where $u$ is the value you got in part \ref{item:4}.
    \item Write a sentence or two interpreting your result.
  \end{enumerate}
\end{problem}

\newpage
\subsection{Hypothesis test type \#1: A one-sided test of parameter value}

\begin{problem}
  The organiziation \textit{Physicians for Human Rights} documented attacks on
  health care facilities during the Syrian civil war. During the period from
  January 2012 through September 30, 2015, the average rate of such attacks by
  Syrian government forces was approximately
  \begin{equation*}
    \theta_0 = 76 \text{ attacks per year}.
  \end{equation*}
  On September 30, 2015, the Russian military began an air campaign in support
  of the Syrian government, substantially increasing aerial operations. From
  that date until December 31, 2016 (a period of $t=1.25$ years), there were
  $X=146$ documented attacks on healthcare facilities by Syrian and Russian
  forces.

  \begin{enumerate}[(a)]
    \setlength{\itemsep}{1em}

    \item Let $\theta$ denote the rate of attacks per year during the time
    period after September 30, 2015. Suppose we wish to assess whether Russian
    intervention coincided with an increase in the rate of attacks on
    healthcare facilities, or whether the rate of such attacks stayed the
    same. State the relevant null and alternative hypotheses.
    \item Let $\widehat{\theta} = X/t$ be the observed rate of attacks in the later
    period, and define the test statistic
    \begin{equation*}
      T = \frac{\widehat{\theta}- \theta_{0}}{\sqrt{\theta_{0}/t}}.
    \end{equation*}
    Make the following two assumptions:
    \begin{itemize}
      \item Assume that if the rate of attacks per year is $\theta$, then the
      number of attacks in a period of $t$ years is Poisson distributed with
      mean $\lambda=t \theta$; i.e., so that $X\sim \text{Poisson}(t\theta).$
      \item Assume that if $Y\sim \text{Poisson}(\lambda)$ and $\lambda>5$ then
      $Y$ is approximately normally distributed with mean $\lambda$ and
      variance $\lambda$.
    \end{itemize}
    Using these two assumptions, show that $T$ has approximately a standard
    normal distribution if $H_{0}$ is true.
    \item Using $T$ as a test statistic, formulate a rejection region of the
    form $\left[ T>k \right]$ at a significance level of $\alpha=.05$. In
    other words, find the value of $k$ such that
    \begin{equation*}
      \P\left[T>k \mid H_{0}\text{ is true}\right] = .05
    \end{equation*}
    Sketch a picture to illustrate your choice of $k$.
    \item Compute the observed value of $T$ and decide whether to reject
    $H_{0}$ or not.
    \item Compute the corresponding $p$-value for your test statistic.
    (Recall, the $p$-value is the probability of observing a value at least as
    large as what you observed, assuming $H_{0}$ is true.)
  \end{enumerate}
\end{problem}

\newpage
\subsection{Hypothesis test type \#3: Test of independence}
\begin{problem}[Test of Independence]
  A 2007 study of $n=186$ autistic Japanese childen examined a possible
  relationship between the onset of autism-related developmental regression
  and exposure to the measles mumps rubella (MMR) vaccine.\footnote{Uchiyama
    et al, \textit{J Autism Dev Disord}, 2007.
    \url{https://doi.org/10.1007/s10803-006-0157-3}} One of their analyeses
  involved the following dataset: \Fl{Table of observed values}
  \begin{center}
    \begin{tabular}{l|l|l|}
      & \textbf{Vaccinated} & \textbf{Not Vaccinated} \\ \hline
      \textbf{Regression}    & $15\phantom{\displaystyle\int}$ & $39$ \\ \hline
      \textbf{No regression} & $45\phantom{\displaystyle\int}$ & $87$ \\ \hline
    \end{tabular}
  \end{center}
  \vspace{1em}
  We'll consider the following hypothesis testing problem:
  \begin{gather*}
    H_{0}: \text{vaccination status and autistic regression are independent}\\
    \text{versus}\\
    H_{a}: \text{there is a relationship between vaccination status and autistic regression}
  \end{gather*}
  \begin{enumerate}[(a)]
    \setlength{\itemsep}{1em}
    \item From the table of observed values, find:
    \begin{itemize}
      \item[]
      \item the total number of children exhibiting autistic regression
      \item the total number of children not exhibiting autistic regression
      \item the total number of vaccinated children
      \item the total number of unvaccinated children
      \item the total number of children
      \item[]  
    \end{itemize}
    \item Let $p$ be the probability that a child exhibits autistic
    regression, and let $q$ be the probability that a child is vaccinated. Use
    your answers from the previous part to estimate $p$ and $q$.
    \item If $H_{0}$ is true, then expected number of children who are both
    vaccinated and exhibit autistic regression is $e_{11} = npq.$ Find
    $e_{11},e_{12},e_{21},$ and $e_{22}$, and arrange them in the following
    table:
    \Fl{Table of expected counts}
    \begin{center}
      \begin{tabular}{l|l|l|}
        & \textbf{Vaccinated} & \textbf{Not Vaccinated} \\ \hline
        \textbf{Regression}    & $e_{11}=\phantom{\displaystyle\int} $ & $e_{12}=$ \\ \hline
        \textbf{No regression} & $e_{21}=\phantom{\displaystyle\int}$ & $e_{22}=$ \\ \hline
      \end{tabular}
      \vspace{1em}
    \end{center}
    \item Compute the chi-squared statistic
    \begin{equation*}
      \chi^{2} = \sum_{\text{all }i,j} \frac{\left( o_{ij}-e_{ij} \right)^{2} }{e_{ij}} 
    \end{equation*}
    where $o_{11},o_{12},o_{21},$ and $o_{22}$ are the observed values. 
    \item Set up and evaluate an integral to compute a $p$-value for your
    $\chi^{2}$ statistic.
    \item Write a sentence or two interpreting your result.
  \end{enumerate}
\end{problem}


\newpage
\begin{problem}[Simpson's Paradox: The Perils of Aggregation]

  In this problem we compare the performance of two medical centers, Hospital
  A and Hospital B. We have the following survival data for surgery patients from
  these two hospitals from the last 6 weeks:
  \begin{center}
    \vspace{1em}
    \begin{tabular}{lcc}
      \hline
      \textbf{Outcome} & \textbf{Hospital A} & \textbf{Hospital B} \\
      \hline
      Died     &   63     & 16       \\
      Survived & 2037       & 784       \\
      \hline
    \end{tabular}
    \vspace{2em}
  \end{center}
  \begin{enumerate}[(a)]
    \setlength{\itemsep}{1em}
    \item \label{item:1} Compare the survival rates of the two hospitals: what percentage of
    patients at Hospital A survive? What about Hospital B?
    \item Is there a relationship between patient survival and which hospital
    they go to? Formulate and perform a chi-squared test of independence. What
    $p$-value do you get? In one or two sentences, state your null hypothesis
    and draw a conclusion about it.
    \item We now consider a third factor: the condition of the patient when
    they are admitted for surgery. Patients are either classified as ``poor
    condition'' or ``good condition''. Here's the more detailed data:

    \begin{center}
      \begin{tabular}{cc}
        \begin{tabular}{lcc}
          \multicolumn{3}{c}{\textbf{Good Condition}} \\ \hline
          \textbf{Outcome} & \textbf{Hospital A} & \textbf{Hospital B} \\ \hline
          Died     &    6    & 8       \\
          Survived & 594       & 592       \\ \hline
        \end{tabular}
                       \quad \quad \quad \quad    &
                             \begin{tabular}{lcc}
                               \multicolumn{3}{c}{\textbf{Poor Condition}} \\ \hline
                               \textbf{Outcome} & \textbf{Hospital A} & \textbf{Hospital B} \\ \hline
                               Died     &  57      & 8       \\
                               Survived & 1443       & 192       \\ \hline
                             \end{tabular}
      \end{tabular}
    \end{center}
    What fraction of Hospital A's patients were classified as good
    condition? What about Hospital B?
    \item \label{item:2} Out of the patients in \textbf{good condition} who went to hospital A,
    what fraction survived? What about hostpial B? Compare these two survival
    rates.
    \item \label{item:3} Out of the patients in \textbf{bad condition} who went to hospital A,
    what fraction survived? What about hostpial B? Compare these two survival
    rates.
    \item Compare your answers to parts
    \ref{item:1}~\ref{item:2}~\ref{item:3}. What you are observing is called
    \demph{Simpson's Paradox}. How can this discrepancy be
    explained? If you are facing surgery, should you go to Hospital A or
    Hospital B?
    \item Fill in the following contingency table:
    \begin{center}
      \vspace{1em}
      \begin{tabular}{|l|c|c|}
        \hline
        \textbf{Outcome} & \textbf{Good Condition} & \textbf{Bad Condition} \\
        \hline
        Died     &    $\phantom{\displaystyle\int}$    &        \\
        \hline
        Survived &        &        $\phantom{\displaystyle\int}$\\ 
        \hline
      \end{tabular}
    \end{center}
    \item Is there a relationship between patient survival and the condition
    that they are admitted for surgery? Formulate and perform a chi-squared
    test of independence and report a $p$-value.
  \end{enumerate}
\end{problem}




\newpage
\begin{problem}[Another example of Simpson's Paradox]
  The influence of race on imposition of the death penalty for murder has been
  much studied and contested in the courts. The following three-way table
  classifies 326 cases in which the defendant was conficted of
  murder.\footnote{Data from M. Radelet, \textit{American Sociological Review}, 1981.
    \url{https://doi.org/10.2307/2095088}} The three variables are the
  defendant's race, the victim's race, and whether or not hte defendant was
  sentenced to death.

  \begin{center}
    \begin{tabular}{cc}
      \begin{tabular}{lcc}
        \multicolumn{3}{c}{\textbf{White Defendants}} \\ \hline
        \textbf{Victim Race} & \textbf{Yes} & \textbf{No} \\ \hline
        White & 19 & 132 \\
        Black & 0  & 9   \\ \hline
      \end{tabular}
                            \quad \quad \quad \quad &
                               \begin{tabular}{lcc}
                                 \multicolumn{3}{c}{\textbf{Black Defendants}} \\ \hline
                                 \textbf{Victim Race} & \textbf{Yes} & \textbf{No} \\ \hline
                                 White & 11 & 52  \\
                                 Black & 6  & 97  \\ \hline
                               \end{tabular}
    \end{tabular}
  \end{center}
  \vspace{3em}
  \begin{enumerate}[(a)]
    \setlength{\itemsep}{7em}
    \item From these data make a 2-way table of defendant's race by death
    penalty. Does it look like there is a relationship between these two
    things? 
    \item Suppose you performed a chi-squared test of independence and for
    the table you just made. You get $p=.64$. How do you interpret this
    result?
    \item Show that Simpson's paradox holds: specifically, that
    \begin{itemize}
      \item a higher percent of white defendants are sentenced to death
      overall
      \item but for both black and white victims, a higher percent of black
      defendants are sentenced to death.
    \end{itemize}
    \item Basing your reasoning on the data, explain why the paradox holds
    in language that a judge could understand.
  \end{enumerate}


 
\end{problem}


% \newpage
% \begin{problem}[One True Love]
  
%   A poll is taken of $n=2625$ individuals; people were survedy on some
%   demographic information and also asked whether or not they believed in ``the
%   one true love''.

%   In this problem, we test whether or not there is a relationshp between the
%   following two factors
%   \begin{itemize}
%     \item Factor 1: education level
%     \item Factor 2: sentiment about true love
%   \end{itemize}

%   We can construct a two-way contingency table for these two factors, as follows:
%   \begin{center}
%     \begin{tabular}{l|r|r|r|r}
%       \textbf{Response} & \textbf{HS} & \textbf{Some College} & \textbf{College} & \textbf{Total} \\
%       \hline
%       Agree       & 363 & 176 & 196 & 735 \\
%       Disagree    & 557 & 466 & 789 & 1812 \\
%       Don't Know  &  20 &  26 &  32 & 78 \\
%       \hline
%       \textbf{Total} & 940 & 668 & 1017 & 2625 \\
%     \end{tabular}
%   \end{center}
%   This contingency table has $r=3$ rows and $c=3$ columns (plus a row and
%   column for totals). The data counts themselves take the form of a $3\times 3$ matrix
%   \begin{equation*}
%     \begin{matrix}
%       Y_{11} & Y_{12} & Y_{13} & \cdots & Y_{1c}\\
%       Y_{21} & Y_{22} & Y_{23} & \cdots & Y_{2c}\\
%       \vdots &&&&\vdots \\
%       Y_{r1} & Y_{r2} & Y_{r3} & \cdots & Y_{rc}\\
%     \end{matrix}
%   \end{equation*}
%   where $Y_{ij}$ is the number of samples for which factor $1$ is $i$ and
%   factor 2 is $j$. We've also added row and column totals in the margins. We
%   can construct a coresponding table of ``estimated expected values'' of the form

%   \begin{equation*}
%     \begin{matrix}
%       E_{11} & E_{12} & E_{13} & \cdots & E_{1c}\\
%       E_{21} & E_{22} & E_{23} & \cdots & E_{2c}\\
%       \vdots &&&&\vdots \\
%       E_{r1} & E_{r2} & E_{r3} & \cdots & E_{rc}\\
%     \end{matrix}
%   \end{equation*}
%   by the formula
%   \begin{equation}\label{eq:1}
%     E_{ij} = n\times \frac{(\text{sum of row }i)}{n} \times \frac{(\text{sum of column }j)}{n}
%   \end{equation}
%   which gives us a new table of estimated expected values.

%   \Fl{Fact:} For contingency tables with $r$ rows and $c$ columns, the chi-squared
%   statistic
%   \begin{equation*}
%     \chi^{2} := \sum_{i=1}^{r}\sum_{j=1}^{c} \frac{\left( Y_{ij}-E_{ij} \right)^{2} }{E_{ij}}  
%   \end{equation*}
%   has chi-squared distribution with $(r-1)\times(c-1)$ degrees of freedom.

%   \fl{Questions:}
%   \begin{enumerate}[(a)]
%     \setlength{\itemsep}{3em}
%     \item Explain why the formula in \Cref{eq:1} holds.
%     \item What is the degrees of freedom for the chi-squared statistic in this
%     problem? Write down the formula for the pdf of the chi-squared statistic.
%     \item Do the $\chi^{2}$ test of independence (i.e., compute a table of
%     expected values, find $\chi^{2}$, and use the known distribution of the
%     statistic to compute a $p$-value). 
%     \item Intepret your result.
%   \end{enumerate}
% \end{problem}
\end{document}